\section{関連研究}
\label{sec:related_work}

前章では，線形倒立振子モデル（LIPM）を基礎として，
二足歩行ロボットにおける歩行制御の枠組みを整理した．
LIPM は，重心高さが一定であるという仮定のもとで，
ZMP と重心運動の関係を簡潔に記述できる点において，
歩行生成および安定化制御に大きく貢献してきた．

一方で，実環境における歩行を考えた場合，
重心高さが常に一定であるという仮定は必ずしも成立しない．
特に，段差や不整地，あるいは床の変形や沈み込みを伴う環境では，
歩行中に重心高さが時間的に変動することが避けられず，
このような状況では，従来の LIPM に基づく制御手法の適用に
課題が生じることが指摘されている．

このような背景から，
重心高さを一定と仮定するのではなく，
重心高さの変動を陽に考慮した歩行モデルや制御手法に関する研究が
数多く提案されてきた．
これらの研究では，
重心高さの調整を自由度として導入することで，
歩行安定性の向上や外乱へのロバスト性の改善，
さらにはエネルギー効率の向上を図ることが試みられている．

本章では，
重心高さが可変であることを前提とした二足歩行制御に関する
代表的な関連研究を取り上げ，
それぞれのアプローチと特徴について整理する．


\section{重心高さを動的に変化させるモデル（Variable-height Inverted Pendulum）}

この問題に対処するため，重心高さを可変とする倒立振子モデルが提案されてきた．
代表的なものとして，Caron らによる
Variable Height Inverted Pendulum（VHIP）モデルが挙げられる
\cite{caron2017,caron2019}．
VHIPでは，鉛直方向の加速度を制御入力として導入し，
重心高さの変化を積極的に利用することで，
段差地形や外乱に対する安定性向上を図っている．

さらに，VHIPを基に，Capture Point（CP）や
Divergent Component of Motion（DCM）の概念を
可変重心高さ環境へ拡張した研究も行われている．
CP や DCM は，重心運動を安定成分と不安定成分に分解し，
ロボットのバランス状態を代表点として表現する指標であり，
LIPM に基づく歩行制御において広く用いられてきた．
これらの概念を重心高さ可変モデルへ拡張することで，
捕捉可能性や安定条件を解析的に評価する手法が提案されている
\cite{liu2020}．

これらの研究は理論的に強固であり，
重心高さ変化を伴う動的バランス制御に対して高い表現力を有する．
一方で，VHIPは運動方程式が非線形となるため，
制御則設計や最適化問題の定式化が複雑化する．
その結果，計算コストの増加や実装負担が大きくなり，
既存のLIPMベース歩行制御器へ容易に統合できないという課題が指摘されている．

\section{LIPMを基盤とした鉛直方向拡張と補助入力の導入}

可変重心高さモデルとは別に，
LIPMの枠組みを維持したまま，
鉛直方向の運動や角運動量を補助的に導入する研究も行われている．
Englsberger らは，
LIPMに基づく Capture Point 追従制御に，
上体角運動量および鉛直方向の重心運動を統合することで，
外乱耐性を向上させる手法を提案している \cite{englsberger2015}．

また，Tedrake らや Pratt らの研究では，
Prismatic Inverted Pendulum Model（PIPM）を用い，
重心が可変高さの軌道上を移動するモデルとして歩行を表現している
\cite{dai2014}．
これらの研究では，
LIPM由来の低次元モデルを基盤としつつ，
高さ変化を幾何学的制約として扱う点に特徴がある．

これらのアプローチは，
LIPMの計算効率や構造的単純さを保ちながら，
鉛直方向の自由度を取り込む点で有効である．
しかし，床沈み込みのように，
支持脚高さが連続的かつ予測困難に変動する状況を
直接対象とした研究は多くない．

\section{コンプライアント地面を考慮した Dual-SLIP に基づくモデルベース歩行制御}
\label{subsec:related_dualslip_compliant}

軟弱地面上歩行に対するモデルベース研究として，Karakasis らは 3 次元 Dual-SLIP（Dual Spring-Loaded Inverted Pendulum）モデルを拡張し，地面のコンプライアンス（剛性および減衰）を明示的に取り込んだ上で，ロバストな動的歩行を実現する制御手法を提案している \cite{karakasis2022med}．本研究では「沈み込み床」を Kelvin--Voigt 型の粘弾性として扱っているが，同様に Karakasis らの枠組みも，地面を剛体ではなく可変剛性・減衰を持つ接触環境としてモデル化する点に特徴がある \cite{karakasis2022med}．

同論文では，片脚支持期における片側の地面剛性が 1 歩だけ低下するような摂動（one-step unilateral stiffness perturbation）を想定し，この摂動下で歩容が破綻しないことをロバスト性の観点から評価している \cite{karakasis2022med}．さらに，提案手法は生体模倣的な方針に基づき，歩行中に脚の剛性を調整することでコンプライアント地面への適応性を高める構成となっている \cite{karakasis2022med}．ベースラインとして LQR 制御との比較が行われ，LQR は中程度以上の地面剛性でのみ安定であるのに対し，提案手法はより低い剛性条件まで安定歩行を維持できることが報告されている \cite{karakasis2022med}．

この研究は，軟弱地面を「モデルのパラメータ（剛性・減衰）」として明示的に取り込み，その不確かさや急激な変化に対しても歩行を破綻させない制御設計を行うという点で，軟弱地面適応に関する代表的なモデルベースアプローチである．一方で，Karakasis らの手法は，3 次元 Dual-SLIP モデルを基盤とし，脚の弾性特性を調整することで地面コンプライアンスへの適応を実現している点に特徴がある．

これに対し，本研究では，歩行制御の基盤として広く用いられている線形倒立振子モデル（LIPM）に着目し，その基本構造を維持したまま，沈み込みに起因して生じる歩行の不安定化に対する適応性を拡張することを目指す．

したがって，本研究は，Dual-SLIP のように新たな力学モデルを導入するのではなく，既存の LIPM ベース歩行制御器との親和性を保ちつつ，軟弱地面に対する適応性を拡張する点に特徴がある．この点において，脚剛性の調整を主な適応手段とする Karakasis らの枠組みとは異なるアプローチから，沈み込み床への適応を実現する試みとして位置づけられる．

\section{本研究の位置づけ}

以上の関連研究を踏まえると，
本研究は以下の点において既存研究と異なる特徴を有する．

本研究では，
床沈み込みにより生じる支持脚高さの変動を推定し，
その変動成分を重心高さの参照値に反映させる
補償型アプローチを採用する．
この際，VHIPのようにモデル自体を非線形に拡張するのではなく，
LIPMの運動方程式および水平方向の制御構造を維持したまま，
鉛直方向の速度および加速度成分を
目標重心高さに付加する．

この設計により，
既存のLIPMベース歩行制御器へ容易に統合可能であり，
計算コストの増加を抑えつつ，
軟弱地面に固有の鉛直方向の不確実性に対応できる．
したがって，本研究は，
可変重心高さを明示的にモデル化する研究と，
重心高さ一定を前提としたLIPM制御の中間に位置づけられ，
床沈み込みという実環境特有の問題に対する
実装容易な解決策を提供するものである．

